\chapter{Parcours technicien}

\section{Responsabilité}

Le technicien exécute les tâches de terrain qui lui sont assignées. Son accès aux
stations est volontairement en lecture seule pour les données de référence. Il
peut consulter le contexte OCPP et les alertes utiles, mais ne crée ni station,
ni connecteur, ni compte utilisateur.

\manualscreenshot{technician-overview.png}{Tableau de bord du technicien}{fig:user-technician-overview}

\section{Consulter les travaux assignés}

Le tableau de bord met en avant les alertes assignées, travaux actifs, échéances
et éléments en retard. La page \textbf{My interventions} permet de filtrer par
statut et d'ouvrir la fiche de travail.

\manualscreenshot{technician-interventions.png}{Liste des interventions assignées}{fig:user-technician-interventions}

\begin{procedurebox}{Exécuter une intervention}
  \begin{enumerate}[leftmargin=1.6em]
    \item ouvrir l'intervention assignée ;
    \item vérifier la station, le connecteur, la priorité, le SLA et le
    diagnostic initial ;
    \item passer l'intervention en cours au début du travail ;
    \item consigner les contrôles, pièces et actions effectuées ;
    \item joindre les photos \textbf{avant} et \textbf{après} si elles sont
    requises ;
    \item rédiger le résultat et les recommandations ;
    \item soumettre le rapport puis terminer l'intervention.
  \end{enumerate}
\end{procedurebox}

\expected{L'opérateur et l'administrateur retrouvent l'historique, les preuves et
le rapport. L'alerte liée peut ensuite être résolue selon son état réel.}

\section{Consulter le contexte d'une station}

\manualscreenshot{technician-stations.png}{Inventaire des stations visible par le technicien}{fig:user-technician-stations}

La fiche station permet de vérifier :

\begin{itemize}[leftmargin=1.6em]
  \item la dernière communication et le dernier \emph{heartbeat} ;
  \item l'état calculé de la station et des connecteurs ;
  \item les erreurs actives et événements récents ;
  \item la maintenance en cours ;
  \item les documents techniques disponibles.
\end{itemize}

\attention{Les graphiques de télémétrie représentent les événements réellement
collectés. Une absence de mesure n'est pas une valeur nulle : elle doit être
interprétée avec la connectivité et les événements OCPP.}

\section{Diagnostiquer avec le Simulation Lab}

Le technicien peut ouvrir le laboratoire pour une station simulée de son
organisation. Son niveau \textbf{Diagnostic} permet d'observer les impulsions,
envoyer un heartbeat, reproduire un branchement ou un débranchement, injecter un
défaut de connecteur, rétablir le connecteur et lancer les scénarios de test.

La connexion ou déconnexion globale de la station, le redémarrage, le
déverrouillage et le mode maintenance restent désactivés : ces actions relèvent
de l'opérateur ou de l'administrateur.

\begin{procedurebox}{Vérifier un défaut et son rétablissement}
  \begin{enumerate}[leftmargin=1.6em]
    \item choisir la station et le connecteur concernés ;
    \item noter l'état et l'heure du dernier signal ;
    \item utiliser \textbf{Inject fault} dans le contexte de test autorisé ;
    \item vérifier l'impulsion \textbf{StatusNotification}, l'état
    \textbf{Faulted} et l'alerte éventuelle ;
    \item utiliser \textbf{Recover} ;
    \item confirmer le retour à \textbf{Available} et consigner le résultat dans
    l'intervention.
  \end{enumerate}
\end{procedurebox}

\section{Rapport de terrain}

\manualscreenshotcompact{technician-reports.png}{Bureau de rapports de maintenance}{fig:user-technician-reports}

Le rapport doit rester factuel et vérifiable :

\begin{itemize}[leftmargin=1.6em]
  \item symptôme observé et heure ;
  \item diagnostic et contrôles réalisés ;
  \item action corrective et pièces utilisées ;
  \item état final de la station et du connecteur ;
  \item photos ou documents joints ;
  \item risque résiduel et recommandation.
\end{itemize}

Le bureau de rapports permet aussi d'envoyer une synthèse à un opérateur ou à
l'administrateur. Les pièces jointes complémentaires sont privées et soumises à
une nouvelle vérification d'autorisation lors de chaque téléchargement. Les
photos \textbf{avant} et \textbf{après} restent, elles, liées au parcours guidé
de l'intervention.

\securitynote{Masquer les visages, plaques d'immatriculation et informations de
paiement non nécessaires. Ne jamais photographier un secret OCPP ou une clé
d'accès lisible.}
